% !TeX root = ../papers/nc-p3.tex
\section{二平面分解の可逆部における全正則性と相補逆行列}
\label{sec:split-inverse-windows}
四位置で $V_i=X_i\oplus Y_i$、両者二次元、重みを $0,1$ とする。
全体重み2、各辺の色階数 $(1,4,1)$ とし、一文字全階数を仮定する。
各辺の純次数関係を $r_i\in X_i\otimes X_{i+1}$,
$s_i\in Y_i\otimes Y_{i+1}$ とし、混合次数の関係四平面を $M_i$ とする。
\begin{theorem}\label{thm:split-inverse-regularity}
全ての $r_i,s_i$ が行列として可逆であることと、
全位置で $M_i\to Y_i\otimes X_{i+1}$ が同型であることは同値である。
この場合 $M_i\to X_i\otimes Y_{i+1}$ も全て同型で、
これらの窓は係数への追加条件なしに全て標準Koszul左右AS正則である。
\end{theorem}
\begin{proof}
全体の二平面別の六成分のうち $XXYY, XYYX, YYXX,YXXY$ は、
色0と2の切断階数が1なので、各々対応する $r_i\otimes s_{i+2}$ の位置置換・非零倍である。
混合対混合の $8\times8$ 切断の二交差ブロックはこれらの二次行列二つの
Kronecker積の置換である。各因子が可逆なら両交差ブロックが可逆で、切断階数が4なので
列空間・行空間の両半分への射影は全て同型である。

逆に一方の射影だけを全位置で仮定する。二次双対では双対文字の $XY$ を $YX$ に
常に書き直せる（元の $M_i$ を $YX$ 上のグラフとして直交補を取る）。
双対の任意の語は従って全 $Y$ の後に全 $X$ が来る語で張られる。
もし $r_i$ が特異なら、その右核に非零の $x\in X_{i+1}^*$ があり、$X_i^*x=0$ である。
位置 $i+1$ の $x$ に終わる四文字で最上形式が評価できるのは $Y$ が二つのときだけである。
その前の三文字を $YYX$ に並べ替えると全て零で、最上テンソルの一文字全階数に反する。
この評価は二次双対の関係を消すので、まだKoszul性を要しない。$s_i$ が特異な場合は
左核の $y$ から始めて後続三文字を $YXX$ に並べれば同じ矛盾を得る。

元の代数で $YX$ は $XY$ に書き直せる。二本の非退化二文字関係列はそれぞれ、
整数位置に沿って基底を順に選ぶと通常の交代関係になる。高位文字を低位文字より大きくした
位置別左辞書順を使えば混合四式と両半分の各一式はすべて降順対を先頭語に持つ。
従って長さ3での正規候補は20語である。最上テンソルの一文字縮約4本は独立な三次重なりを与え、
実際の次元も20以上である。全三文字の重なりが解消し、Diamond補題で全次数の二次PBW基底を得る。
四位置の有限Frobenius完全対合は、Koszul双対次元が $1,4,6,4,1$ となった後に
元の切断階数と最上形式から従う。
\end{proof}

\begin{theorem}\label{thm:split-inverse-transporter}
この可逆部で、順序付き二平面分解を保つ添字同型類は次のデータで全て与えられる。
四つの二次元空間 $U_i$ 上の射影テンソル $P$ で、$01|23$ と $12|30$ の両切断が可逆なものを取る。
それぞれの切断で逆行列の転置を四つの双対空間に戻した射影テンソルを $I_0(P),I_1(P)$ とする。
$D=\prod\operatorname{PGL}(U_i)$ とし、二次元空間の射影線とその双対は交代形式で同一視する。
$\tau=(02)(13)$ と書き、$U_i$ を固定した二次元模型に移す。
\[
  H=\operatorname{Stab}_D(P),\qquad
  \mathcal B(P)=\{v\in \tau D:\ v I_1(P)=I_0(P)\}.
\]
$P$ を $D$ 同値まで取り、$\mathcal B(P)$ が空なら捨てる。各 $P$ の上の
同型類は $\mathcal B(P)$ の $H$ 共役類である。非空なら $\mathcal B(P)=Hv$、
$vHv^{-1}=H$。これは非半単純な $P$ にも成立する。
\end{theorem}
\begin{proof}
$U=(X_0,Y_1,X_2,Y_3)$ とし、各位置の相補空間を $Z$ と書く。
交代二成分を $P\in\bigotimes U_i$, $Q\in\bigotimes Z_i$ とする。
第一切断の交差成分は非退化な二文字形式の積であり、
$\pi_0=(03)(12)$ の四つの対応で射影同型 $L_i^0:U_i^*\to Z_{\pi_0(i)}$ を与える。
第二切断では $\pi_1=(01)(23)$ の四対応 $L_i^1$ となる。
各切断で $P$ のブロックを左上とした行・列順にする。交差ブロックが可逆で全階数が4なので
$P,Q$ の当該ブロックも可逆であり、Schur式 $Q=V P^{-1}U$ が成り立つ（右辺の $U,V$ は二交差ブロック）。
添字を戻すと $[Q]=\pi_j(\bigotimes L_i^j)[I_j(P)]$ である。
\path{computations/verify_split_inverse_windows.py} は両方の添字順を直接検算する。

逆に、二つの射影等式を満たすデータから、六成分の代表倍率を決められる。
$P,Q$ の代表を固定し、各切断で二交差成分の一方を倍率変更してSchur式を厳密に合わせる。
二切断で使う交差成分は別々なので両立する。全解は一つの添字同型類を与える。
実際、高位位置が $i,j$ の六成分の代表の比を $c_{ij}$ とすれば
$c_{01}c_{23}=c_{02}c_{13}=c_{03}c_{12}$ である。これらは $c_{ij}=d_i d_j$ であるための
必要十分条件であり、例えば $d_0^2=c_{01}c_{02}/c_{12}$ から順に取れる。
四位置の高位を $d_i$ 倍すればよい。再構成した窓は純次数の全行列が可逆であり、
各文字はそれらの積で検出されるので一文字全階数も持つ。

第一の四つの $L_i^0$ を相補空間の座標に使うと、残る一組は
$\pi_0^{-1}\pi_1=\tau$ を伴う $v$ となる。
$P$ を変更する残りの $D$ は $v$ に共役で作用する。逆転置写像は射影双対上で
$D$ 同変かつ各切断で対合なので $I_0(P),I_1(P),P$ の安定群は同じ $H$ である。
従って求める移送集合と共役商が得られる。ここで $\tau$ は相補文字の同定から生じ、
窓の始点を移動する操作ではない。
\end{proof}

\begin{corollary}\label{cor:split-finite-unmarking}
一つの窓におけるこの分解の選択は、射影安定群の恒等成分での共役まで有限個である。
$H$ が有限なら順序付き分解を忘れる追加操作は全位置の $X,Y$ の同時交換だけである。
\end{corollary}
\begin{proof}
各分解は重み差1の一径数で決まり、全一径数は安定群の固定した極大トーラスに
共役で入る。そのトーラスの各位置の矢の射影重み差は指標格子の有理空間を張る。
これら有限個の重み差の値が $0,\pm1$ に限られるので、該当する余指標は有限個である。
$H$ が有限なら指定一径数を中心化する群の恒等成分はその一径数だけである。
実際、八つの二次元空間の射影作用は $P$ と四同定から $H$ に入り、残るスカラーで
テンソル直線を保つには六つの高位対の倍率 $d_i d_j$ が全て等しく、全 $d_i$ が等しい。
よって安定群の階数は1で、原始余指標は符号まで一つである。
\end{proof}
